\section{Conclusion}\label{sec:conclusion}

We built a demo bank with four interchangeable transaction-authentication modes
and evaluated them under one threat model with an automated harness, a live
man-in-the-browser demonstration, an online-guessing experiment, and
micro-benchmarks, with every number produced by \texttt{make eval}.

The measurements answer the research questions. \textbf{RQ1:} binding the code to
the transaction stops the recipient and amount tampering that a plain OTP lets
through --- \texttt{hmac} and \texttt{sig} block S2a, S3, S4 and S5, which
\texttt{baseline} does not --- and, by the A1/A2 equivalence, this covers
real-time phishing as well. \textbf{RQ2:} asymmetric signatures add what a
symmetric code cannot --- survival of a database leak (S10) and
third-party--verifiable evidence (S11) --- because the key never leaves the
phone, at a signing cost of about $19\,\mu s$ on the host and under a millisecond
on a phone. \textbf{RQ3:} the trusted display is a separate and necessary layer:
\texttt{sig\_blind}, with identical cryptography but no display, fails exactly the
attacks an attentive user would otherwise catch.

The honest reading is equally important: the guessing limit is only per
transaction, ``verifiable evidence'' is not non-repudiation, the symmetric keys
would need an HSM, and the attentive-user result is an upper bound. The strongest
cryptography does not protect a user who is shown nothing, or who does not look.

Future work follows the limitations: a per-user guessing counter with a
rate-limited creation path, an HSM for the symmetric keys, an out-of-band
enrolment and recovery flow, independent attestation of the signing key to move
from verifiable evidence toward non-repudiation, and a proper user study to
replace the idealised attentive-user model with measured detection rates.
